\subsection*{1.3. 函子 $R\Psi$}

\subsubsection*{1.3.1.}
设 $D$ 是圆盘，$0\in D$，令 $D^*=D-\{0\}$，设 $f:X\longrightarrow D$ 是连续映射，并令 $X_0=f^{-1}(0)$。函子 $R\Psi$ 与 $R\Psi_\eta$ 分别是 $\Psi$ 与 $\Psi_\eta$ 的导出函子：
\[
R\Psi:D^+(X)\longrightarrow D^+(X_0\times\mathcal D),
\]
\[
R\Psi_\eta:D^+(X^*)\longrightarrow D^+(X_0\times\mathcal D^*).
\]
用 $R\Phi$ 表示复合函子 $\Phi\circ R\Psi$：
\[
R\Phi:D^+(X)\longrightarrow D^+(X_0\times\mathcal D^*).
\]
这些函子的性质与 XIII, 2.1 中同名函子的性质完全类似。证明要初等得多，因为在拓扑情形，固有态射或光滑态射（即局部同构于 $\mathbb R^n\times S\longrightarrow S$ 的态射）的基变换定理容易证明。所需的少量改动如下。

\subsubsection*{1.3.2.}
函子 $Rf_!$ 只对局部紧空间之间的分离态射定义，是具有固有支集的正像函子 $f_!$ 的导出函子。

\subsubsection*{1.3.3.}
对固有（且分离）态射 $f:X\longrightarrow D$，基本同构（XIII, (2.1.8.3)）须改写如下。
设 $p:\widetilde D^*\longrightarrow D^*$ 是 $D^*$ 的一个万有覆盖。对 $D$ 中 $0$ 的每个邻域 $U$，令
\[
\widetilde U^*=p^{-1}(U),
\qquad
X_{\widetilde U^*}=X\times_D\widetilde U^*.
\]
对 $K\in\operatorname{Ob}D^+(f^{-1}(D^*))$，有
\begin{equation}\tag{1.3.3.1}
\varinjlim_U\mathbb H^i(X_{\widetilde U^*},K)
\xrightarrow{\ \sim\ }
\mathbb H^i(X_0,R\Psi_{\bar\eta}(K)).
\end{equation}
当 $X^*=f^{-1}(D^*)$ 是 $D^*$ 上的拓扑纤维丛时，此同构尤其有用。当圆盘 $U\subset D$ 缩向 $0$ 时，$X_{\widetilde U^*}\longrightarrow\widetilde U^*$ 的 Leray 谱序列表明：对 $X^*$ 上的每个局部常值层 $\Lambda$ 和每个 $t\in\widetilde U^*$，
\[
H^i(X_{\widetilde U^*},\Lambda)
\xrightarrow{\ \sim\ }H^i(X_t,\Lambda).
\]
因此，对 $t\in D^*$，
\begin{equation}\tag{1.3.3.2}
H^i(X_t,\Lambda)
\xrightarrow{\ \sim\ }
H^i(X_0,R\Psi_t(\Lambda)),
\end{equation}
只要 $X/D$ 固有、$X^*$ 为纤维丛且 $\Lambda$ 局部常值即可。

实际验证 $X^*$ 是 $D^*$ 上的纤维丛时，使用 Thom 的同痕定理，其一种形式陈述如下。关于这个定理的证明，参见 Whitney 和 Mather；其任何推论均不会用于后续各讲。

